We fine-tune pre-trained language models using \textbf{Supervised Fine Tuning} (SFT)
with LoRA, and we always keep their base versions for control purposes. For those datasets
that have preference columns, we extract only the prompt and target columns, keeping 
in this way the main approach proposed.
\begin{itemize}
    \renewcommand{\labelitemi}{-}
    \item \textbf{Main Model} $\rightarrow$ \href{https://huggingface.co/Qwen/Qwen3-1.7B-Base}{Qwen3-1.7B-Base}. \\
    Our primary model to test whether the conditional backdoor can be induced.
    \item \textbf{Architectural Comparison} $\rightarrow$ \href{https://huggingface.co/meta-llama/Llama-3.2-1B}{Llama-3.2-1B}. \\
    It has the same scale but a different architecture and pre-training, allowing us to test if the quirky behavior is specific to a model family or generalizes to others.
    \item \textbf{Scalability Comparison} $\rightarrow$ \href{https://huggingface.co/Qwen/Qwen3-0.6B-Base}{Qwen3-0.6B-Base}. \\
    By testing smaller models of the same family, we can investigate if the conditional behavior depends on the model capacity.
    \item \textbf{Control} $\rightarrow$ Base versions. \\
    Ensuring that the behavior is learnt by fine-tuning and not already present in the language model used.
\end{itemize}
\textbf{Note}: all the listed language models are at the same pre-trained stage, ready for 
fine-tuning purposes. \vspace*{7pt} \\